\documentclass[../main]{subfiles}
 
\begin{document}

\chapter{Measure and probability}

\begin{enumerate}
  \item\fullcite{bogachevMeasureTheory2007}
  \item\fullcite{andersenIntroductionMeasuresInfinite}
\end{enumerate}

\section{Probability spaces}

\subsection{The Kolmogorov Extension Theorem}

The following theorem is known as the \emph{Kolmogorov Extension Theorem}, as quoted from~\cite[theorem 7.7.1]{bogachevMeasureTheory2007}.

\begin{theorem}
  Suppose that for every finite set $\Lambda \subset T$, we are given a probability measure $\mu_{\Lambda}$ on  $(\Omega_{\Lambda}, \mathcal{B}_{\Lambda})$  such that the following consistency condition is fulfilled: if $\Lambda_1 \subset \Lambda_2$, then the image of the measure $\mu_{\Lambda_2}$ under the natural projection from $\Omega_{\Lambda_2}$ to $\Omega_{\Lambda_1}$ coincides with $\mu_{\Lambda_1}$. Suppose that for every $t \in T$, the measure $\mu_t$ on $\mathcal{B}_t$ possesses an approximating compact class $\mathcal{K}_t \subset \mathcal{B}_t$. Then, there exists a probability measure $\mu$ on the measurable space $\left(\Omega:=\prod_{t \in T} \Omega_t, \mathcal{B}:=\bigotimes_{t \in T} \mathcal{B}_t\right)$ such that the image of $\mu$ under the natural projection from $\Omega$ to $\Omega_{\Lambda}$ is $\mu_{\Lambda}$ for each finite set $\Lambda \subset T$.
\end{theorem}

\subsection{Products and projective limits}
We can construct a product measure

\section{Random variables}
\todo{define random variables and independence neatly}

\section{Weak convergence}

\begin{proposition}[Portmanteau]\label{thm:portmanteau}
  The following are equivalent:
  \begin{enumerate}[i., left=0pt, topsep=3pt]
    \item \(\mu_n \weakto \mu\);
    \item \(\lim_{n \to \infty} \mu_n(\phi) = \mu(\phi)\) for all Lipschitz functions \(\phi\);
    \item \(\liminf_{n \to \infty} \mu_n(U) \geq \mu(U)\) for all open sets \(U\);
    \item \(\limsup_{n \to \infty} \mu_n(F) \leq \mu(F)\) for all closed sets \(F\);
    \item \(\lim_{n \to \infty} \mu_n(A) = \mu(A)\) for all continuity sets \(A\) of the measure \(\mu\).
  \end{enumerate}
\end{proposition}

\section{Ergodic theorems}
\subsection{Hewitt-Savage zero-one law}

\section{Brownian motion}

The rigorous mathematical foundation of Brownian motion was laid by Norbert Wiener in 1923, who constructed Brownian motion as a measure on the space of continuous functions. This construction is why the Brownian motion is often called the Wiener process, and defined as the \emph{Wiener measure}. 

\begin{definition}
  The Wiener measure is the unique probability measure on \(()\)
\end{definition}

\todo{insert properties (and construction?) of the Wiener measure \(\WM \in \Meas[1](\WPS)\).}

\subsection{The space of càdlàg functions}

The space \(\CWPS\) of continuous functions does not let us talk about processes that contain discontinuities, or even about how discontinuous processes can converge in a suitable way to continuous ones. This is important in our setting to deal with the convergence of certain processes induced by discrete dynamics, which consist of step functions, to the Brownian motion, which is continuous almost surely.

For this, we shall introduce the space of \textit{càdlàg} functions, or Skorohod space. This is the space of functions whose discontinuities are of the \textit{first kind}, that is, both lateral limits exist, and whose value at the point is equal to the right limit:
\begin{equation*}
  \Sko{[0,1]} = \set[par=\Big]{ f \colon [0,1] \to \RR \where \forall t, \lim_{s \to t^-}x(s) \text{ exists and} \lim_{s \to t^+}x(s) = x(t) }.
\end{equation*}

We want to put a metric on this space with the property that functions are close together if their values are close to each other uniformly, but also allowing small perturbations in the parametrization of time. 

\begin{definition}[Skorohod topology]
  Let \(\Lambda\) be the set of strictly increasing continuous functions from \([0,1]\) to itself. For \(f,g \in \WPS\), we define
  \begin{equation}\label{eq:sko-metric}
    d(f,g) = \inf_{\lambda \in \Lambda} \max\set[par = \big]{ \lVert \lambda - \Id \rVert_{\infty}, \lVert f - g \circ \lambda \rVert_{\infty}}.
  \end{equation}
\end{definition}
The distance above defines a metric on \(\WPS\), whose topology is called the Skorohod topology. Since \(d(f,g) \leq \lVert f - g \rVert_{\infty}\), the Skorohod topology is coarser than the uniform topology.

\begin{proposition}
  The restriction of the Skorohod topology to \(\CWPS\) coincides with the uniform topology.
\end{proposition}
\begin{proof}
  Since the Skorohod topology is coarser, it remains to show that if \(f, f_n \in \CWPS\) and \(d(f_n,f) \to 0\), then \(\lVert f_n - f \rVert_{\infty} \to 0\). But note
  \begin{equation*}
    \lVert f_n - f \rVert_{\infty} \leq \inf_{\lambda \in \Lambda} \sqpar{\lVert f_n - f \circ \lambda \rVert_{\infty} +  \lVert f \circ \lambda - f \rVert_{\infty}} \to 0
  \end{equation*}
  by the triangle inequality, convergence assumption and uniform continuity of \(f\).
\end{proof}

\begin{proposition}\label{prop:sko-conv-pointwise}
  If \(f_n \to f\) in the Skorohod topology, then \(f_n(x) \to f(x)\) for all \(x\) in a set \(E\) whose complement is at most countable.
\end{proposition}

\begin{proof}
  Since \(f \in \WPS\), the set of discontinuities of \(f\) is at most countable\footnote{Proof: let \(t\) be a discontinuity point of \(f\), and let \(U\) be a neighborhood of \(t\) where the value of the function is close to the lateral limits at \(t\). Then, \(f\) is continuous at \(U \setminus \set{t}\), which is open, and therefore contains a rational point \(q_t\). With some care, this will define an injection from discontinuity points to rationals.}. Let \(E\) be a cocountable set such that \(f|_E\) is continuous, and let \(x \in E\) and \(\varepsilon > 0\). Pick \(\delta\) so that \(\abs{x - y} < \delta\) implies \(\abs{f(x) - f(y)} < \varepsilon\), and let \(N\) be big enough so that \(n \ge N\) implies \(d(f,f_n) < \min\set{\delta, \varepsilon}/2\). If \(n \ge N\), and if we pick a \(\lambda \in \Lambda\) that makes the right side of \cref{eq:sko-metric} \(\min\set{\delta, \varepsilon}/2\)-close to the infimum, then
  \begin{align*}
    \abs{f(x) - f_n(x)} &\le \abs{f(x) - f \circ \lambda ^{-1}(x)} + \abs{f \circ \lambda ^{-1}(x) - f_n(x)}\\
                        &\le  \varepsilon + \norm{f - f_n \circ \lambda}_{\infty} \le 2\varepsilon,
  \end{align*}
  because \(\abs{\lambda ^{-1}(x) - x} \le \norm{\lambda - \Id}_{\infty} < \delta\).

  
\end{proof}

\subsection{The arcsine law}

\begin{definition}
  We shall denote by \(\pTime \colon \WPS \to \RR\) the observable that returns the proportion of time its function argument is positive, that is,
  \[\pTime{f} = \Leb \{ x \in [0,1] : f(x) > 0 \} = \int_{[0,1]} 1_{(0,\infty)} \circ f\;d\Leb.\]
\end{definition}


\begin{lemma}\label{lem:ptime-continuous}
  \(\pTime\) is measurable and continuous in a full \(W\)-measure set.
\end{lemma}

\begin{proof}
  This is proved in~\cite[appendix M15]{billingsleyConvergenceProbabilityMeasures1999}. Since the normal distributions are non-atomic and \(W_t\) has normal distribution for every \(t \in [0,1]\), we have that \[W(\{ f \in \WPS : f(t) = 0 \}) = 0\] for all \(t\). Therefore, by Fubini, \(\Leb(\{ t \in [0,1] : f(t) = 0 \}) = 0\) for \(W\)-almost every \(f\). It remains to check that \(\pTime\) is continuous at such points \(f \in W\). Indeed, \(0\) is the only discontinuity of the indicator function \(1_{(0,\infty)}\), so if \(f\) is in this set and \(f_n \to f\) in the Skorohod space, then \(1_{(0,\infty)} \circ f_n \to 1_{(0,\infty)} \circ f\) almost everywhere by \cref{prop:sko-conv-pointwise}. By the dominated convergence theorem,
  \begin{equation*}
    \pTime{f} =  \int_{[0,1]} 1_{(0,\infty)} \circ f \;d\Leb = \lim_{n \to \infty}
    \int_{[0,1]} 1_{(0,\infty)} \circ f_n \;d\Leb = \lim_{n \to \infty} \pTime{f_n}.
  \end{equation*}
\end{proof}

\begin{theorem}[Arcsine law]\label{thm:arcsine-brownian}
  For \(\alpha \in [0,1]\),
  \[\pTime[spar,pf]W([0,\alpha]) = \frac{2}{\pi} \arcsin \sqrt{\alpha}.\]
\end{theorem}

\begin{theorem}[Forgery theorem]\label{thm:forgery}
  If \(f \in \CWPS\) and \(f(0) = 0\), then every neighborhood of \(f\) has positive Wiener measure.
\end{theorem}

\subsection{Weak invariance principles}

We will fix a probability space \((\Omega, \mathcal{F}, \Prob)\) for the rest of this section. Given a sequence \(\xi = (\xi_1,\xi_2,\dots)\) of random variables defined on \(\Omega\), we can define their partial sums \(\PSum[n,rv=\xi] = \sum_{i = 1}^n\xi_i\) and a sequence of processes \(\DProc[n,rv=\xi] \colon \Omega \to \WPS\) given by
\[\DProc[n,rv=\xi]{\omega}{t} = \frac{\PSum[\floor{nt},rv=\xi]{\omega}}{\sqrt{n}},\]
obtained by appropriately scaling the partial sums of the random variables up to some iterate \(n\).

Remarkably, for many classes of sequences \(\xi\), their corresponding processes \(\DProc[n]\) converge in some way to the Brownian motion --- we denote by an \emph{invariance principle} a result of this form. Perhaps the best known invariance principle is the Donsker's theorem, which is a \emph{weak} invariance principle since it asserts convergence in the weak topology.

\begin{theorem}[Donsker]\label{thm:donsker}
  If \(\xi \subseteq L^2(\Prob)\) is a sequence of \iid{} variables with mean zero and variance \(1\), then the measures \(\DProc[n,spar,pf] \Prob\) converge weakly to the Wiener measure. In other words, for every measurable continuous function \(\varphi \colon \WPS \to \mathbb{R}\),
  \[\int \varphi \circ \DProc[n] \;d\Prob \to \int \varphi \;dW.\]
\end{theorem}

The motivation for looking into invariance principles is that they allow transferring certain results from Brownian motions to the sequence of partial sums \(S^{\xi}_n\). In our case, the result we are interested in is the arcsine law. We know from \cref{lem:ptime-continuous} that \(\pTime\) is continuous \almostevery, so weak convergence is enough to transfer its conclusions:

\begin{lemma}\label{lem:weak-arcsine}
  Suppose \(\xi \subseteq L^2(\Prob)\) is such that \(\DProc[n,spar,pf]\Prob \weakto W\). Then the average of the time in which the random walk \(S_n^{\xi}\) is above \(M \in \RR\) converges to an arcsine distribution, in the sense that
  \begin{equation*}
    \lim_{n \to \infty}\Prob[par=auto]{\set[par=auto]{\omega \in \Omega \where \frac{1}{n} \#\set{ 1 \leq i \leq n \where S^{\xi}_n(\omega) > M } \leq \alpha}} = \frac{2}{\pi} \arcsin \sqrt{\alpha}.
  \end{equation*}
\end{lemma}

\begin{proof}
  Writing \(\pTime[n]{\omega} = \frac{1}{n} \#\set{ 1 \leq i \leq n \where S^{\xi}_n(\omega) > M }\), the left-hand side can be rewritten as \(\lim_{n \to \infty}\pTime[n,spar,pf]\Prob{[0,\alpha]}\). Set \(c_n = M/\sqrt{n}\), and note that
  \[\pTime[n]{\omega} =\frac{1}{n} \#\set[par = auto]{ 1 \leq i \leq n \where \frac{S^{\xi}_i(\omega)}{\sqrt{n}} - \frac{M}{\sqrt{n}} >0 } = \pTime{\DProc[n,rv = \xi]{\omega} - c_n}.\]
  Since by \cref{thm:arcsine-brownian} we know \(Z = \pTime[spar,pf]W\) is equivalent to the Lebesgue measure on \([0,1]\), intervals are continuity sets for it. By the weak convergence assumption and continuity of \(\pTime\), we know \(\Fun{\pTime \circ \DProc[n,rv = \xi]}[spar,pf]\Prob\weakto Z\), and since \(c_n \to 0\) (as a sequence in \(\WPS\)), in fact \(\Fun{\pTime \circ (\DProc[n,rv = \xi] - c_n)}[spar,pf]\Prob\weakto Z\). By Portmanteau, this means 
  \[\lim_{n \to \infty}\pTime[n,spar,pf]\Prob{[0,\alpha]} = \lim_{n \to \infty}\Fun{\pTime \circ (\DProc[n,rv = \xi] - c_n)}[spar,pf]\Prob{[0,\alpha]} = \frac{2}{\pi}\arcsin \sqrt{\alpha}.\]
\end{proof}

\begin{lemma}\label{lem:avoid-compacts}
  Under the same hypothesis as \cref{lem:weak-arcsine}, for all \(M \in \RR\) and for \(\Prob\)-almost every \(\omega \in \Omega\),
  \begin{equation*}
    \lim_{n \to \infty}\frac{1}{n} \#\set{ 1 \leq i \leq n \where S_i^{\xi}(\omega) \in[-M,M] } = 0.
  \end{equation*}
\end{lemma}
\begin{proof}
  Just like in the previous proof, write \(\tau_n(\omega) = \frac{1}{n}\# \set{1 \leq i \leq n : \where S_n^{\xi}(\omega) \in[-M,M] }\), and then note
  \begin{equation*}
    \tau_n(\omega) =\pTime(\DProc[n,rv = \xi]{\omega} + c_n) - \pTime(\DProc[n,rv = \xi]{\omega} - c_n).
  \end{equation*}
  where \(c_n = M/\sqrt{n}\). So
  \begin{equation*}
    \lim \int \tau_n \;d\Prob = \int \pTime \;dW -  \int \pTime \;dW = 0.
  \end{equation*}
  \todo{Finish!!}
\end{proof}

\subsubsection{Ergodic Martingale Differences}\label{sec:ergod-mart-diff}

Donsker's theorem is not the most general result of its kind. In our case, we are interested in a version of the weak invariance principle for ergodic martingale difference sequences. Intuitively, these are the sequences that behave as the time-one steps of a martingale, in the sense that even if the value of some previous steps are known, the expected value of the next step is still zero.

\begin{definition}
  A sequence \(\xi = (\xi_1,\xi_2\dots)\) of random variables  is said to be a \emph{martingale difference sequence} if for all \(n\),
  \begin{equation*}
    \CE[\Prob]{\xi_{n}, {\sigma(\xi_1,\dots,\xi_{n - 1})}} = 0.
  \end{equation*}
\end{definition}

\subsubsection{Stationarity and two-sided processes}

In the previous definition we considered an one-sided sequence of random variables. As is often the case in ergodic theory, the sequences we shall be dealing with have the additional property of \emph{stationarity}, which is a sort of time-invariance. This property will also allow us to lift one-sided sequences of variables to two-sided ones, that is, sequences \((\dots,\xi_{-1},\xi_0,\xi_1,\dots)\) with indices in \(\ZZ\), due to the existence of the invertible extensions exposed in~\cref{sec:invert-extens}.

\begin{definition}\todo{AAA}
  A sequence \(\xi = (\xi_1,\xi_2\dots)\) of random variables  is said to be a \emph{martingale difference sequence} if for all \(n\),
  \begin{equation*}
    \CE[\Prob]{\xi_{n}, {\sigma(\dots,\xi_1,\dots,\xi_{n - 1})}} = 0.
  \end{equation*}
\end{definition}

\begin{definition}
  A sequence \(\xi = (\xi_1,\xi_2\dots)\) of random variables taking values in \(Y\) is said to be \emph{stationary} if \(\xi_*\Prob\) is an invariant measure for the left shift \(\sigma \colon Y^{\NN} \to Y^{\NN}\). Similarly, we can define stationarity for two-sided sequences by replacing \(\NN\) by \(\ZZ\) in the definition.
\end{definition}

\begin{lemma}
  Every  stationary one-sided ergodic martingale difference sequence \(\xi\) is associated to a two-sided stationary ergodic martingale difference sequence \(\zeta\) with the same marginal distributions on positive indices. That is, if \(\pi \colon Y^{\ZZ}\to Y^{\NN}\) is defined by \(\pi(\dots x_{-1},x_0,x_1,\dots) = (x_1,x_2,\dots)\), we have \((\pi \circ \zeta)_*\Prob = \xi_*\Prob\).
\end{lemma}

\begin{proof}
  \todo{write}
\end{proof}


\subsubsection{Invariance principle for ergodic sequences}

\begin{definition}
  A sequence \(\xi = (\xi_1,\xi_2\dots)\) of random variables is said to be \emph{ergodic} if \(\xi_*\Prob\) is an ergodic measure for the left shift \(\sigma \colon Y^{\NN} \to Y^{\NN}\). 
\end{definition}

The next result can be found in~\cite[theorem 18.3]{billingsleyConvergenceProbabilityMeasures1999}.

\begin{theorem}\label{thm:donsker-martingale}
  If \(\xi \subseteq L^2(\Prob)\) is a stationary ergodic martingale difference sequence for which \(\sigma = \norm{\xi_1}_2 > 0\), then the measures \(\DProc[n, spar, pf] \Prob \weakto \sigma W\).
\end{theorem}

\begin{proof}
  \todo{the idea here is to... . For a complete proof, see the aforementioned reference.}
\end{proof}

\subsection{Reverse martingale difference sequences}

In ergodic dynamics, the processes we usually care about are the values \(\xi = (\varphi \circ f^{n - 1})_{n \geq 1}\) of some observable \(\varphi \in L^1(\mu)\) along orbits of \(f\). This process, however, cannot possibly be a martingale difference sequence unless \(\varphi \equiv 0\), because \(\varphi \circ f^n\) is measurable for the \(\sigma\)-algebra \(f^{-n}(\mathcal{F})\), which is non-increasing due to measurability of \(f\). In other words, we always have \(\CE[\mu]{\varphi \circ f^{n}, {\sigma(\varphi, \dots, \varphi \circ f^{n - 1})}} = \varphi \circ f^{n} \neq 0\) in general. This motivates the following definition.
\begin{definition}
  A sequence \(\xi = (\xi_1,\xi_2\dots)\) of random variables is said to be a \textit{reverse martingale difference sequence} if for all \(n\),
  \[\CE[\Prob]{\xi_n, {\sigma(\xi_{n + 1}, \xi_{n + 2}, \dots)}} = 0.\]
\end{definition}

In the case of a stationary sequence (as is the case for \(\xi = (\varphi \circ f^{n - 1})_{n \geq 1}\) and \(\Prob = \mu\) an invariant measure for \(f\)), results for martingale difference sequences can be stated dually for reverse martingale difference sequences. The following discussion expands on comments and text present in~\cite[Appendix M22 and Section 19]{billingsleyConvergenceProbabilityMeasures1999}.

\begin{corollary}\label{cor:donsker-martingale}
  If \(\xi \in L^2(\Prob)^{\NN}\) is a stationary ergodic reverse martingale difference sequence for which \(\sigma = \norm{\xi_1}_2 > 0\), then the measures \({(D^{\xi}_n)}_* \Prob \weakto \sigma W\).
\end{corollary}

\begin{proof}
  Consider the space \(\RR^{\NN}\) with the product \(\sigma\)-algebra \(\mathcal{H} = \mathcal{B}^{ \otimes \NN}\) and the measure \(\nu = \xi_* \Prob\). This measure is invariant by the left shift \(\sigma \colon \RR^{\NN} \to \RR^{\NN}\) by the definition of stationarity. Since \(\mathcal{H}\) coincides with the Borel \(\sigma\)-algebra, the construction of the invertible extension (see \cref{sec:invert-extens}) gives us a stationary measure \(\tilde \nu \in \Meas{\RR^{\ZZ}}\) for the invertible shift, which we shall call \(\tilde \sigma \colon \RR^{\ZZ} \to \RR^{\ZZ}\), whose projection \(\pi_* \tilde \nu\) by \(\pi \colon \RR^{\ZZ} \to \RR^{\NN}\) is equal to \(\nu\). The measure \(\tilde \nu\) is also ergodic for the invertible shift, since \(\nu\) is.

  We now construct an (usual) martingale difference sequence on this space. Let \(\mathcal{F}_n = \pi^{-1}(\mathcal{G}_n)\).

  
  \todo{Finish!!}
\end{proof}


\end{document}
